% Generated draft: requires mathematical and source review.
\documentclass[11pt]{article}
\usepackage[margin=25mm]{geometry}
\usepackage{fontspec}
\setmainfont{Latin Modern Roman}
\usepackage{amsmath,amssymb,mathtools,unicode-math}
\setmathfont{Latin Modern Math}
\usepackage{enumitem,needspace,xurl,hyperref}
\setlength{\parindent}{0pt}
\setlength{\parskip}{5pt}
\setlength{\emergencystretch}{3em}
\newcommand{\N}{\mathbb{N}}
\newcommand{\Z}{\mathbb{Z}}
\newcommand{\Q}{\mathbb{Q}}
\newcommand{\R}{\mathbb{R}}
\newcommand{\C}{\mathbb{C}}
\newcommand{\F}{\mathbb{F}}
\newcommand{\A}{\mathbb{A}}
\newcommand{\PP}{\mathbb{P}}
\newcommand{\E}{\mathbb{E}}
\newcommand{\eps}{\varepsilon}
\newcommand{\dd}{\,\mathrm{d}}
\begin{document}
\section*{UnsolvedMath ID: 800412}
\textbf{Source problem code:} catalogue entry
\subsection{Definitions and mathematical statement}
Write \(K=\mathbb{C}_v\) for the field appearing in the supplied source statement, and let \(f\in K(z)\) be a rational function of degree \(d\geq 2\). The supplied excerpt does not specify the class of fields denoted by \(\mathbb{C}_v\), the ambient analytic space in which components are taken, or the meanings of ``indifferent component'' and ``disk.'' Consequently, the formulation below is conditional on those intended source conventions; these missing definitions must be recovered before the statement is mathematically complete.

Under those conventions, a periodic cycle of components is a finite orbit \((U_0,\ldots,U_{n-1})\) of distinct components under the action induced by \(f\), where \(n\geq 1\) and \(U_i\) maps to \(U_{i+1\bmod n}\). Cyclic relabelings describe the same cycle. Let \(N_K(f)\) be the number, possibly infinite, of such cycles whose components are indifferent and are not disks, interpreted according to the source conventions.

\textbf{Question 35 (Benedetto, DeMarco, Ingram, Jones et al., 2018).} Can \(N_K(f)\) be infinite? Equivalently, does the finiteness assertion
\[
\forall f\in K(z),\qquad \deg(f)\geq 2\ \Longrightarrow\ N_K(f)<\infty
\]
hold for each field \(K=\mathbb{C}_v\) allowed by the source?

If finiteness holds, is the number bounded in terms of the degree? For a fixed admissible field \(K\), this asks whether
\[
\forall d\in\mathbb{Z}_{\geq 2}\ \exists B_K(d)\in\mathbb{Z}_{\geq 0}\ \forall f\in K(z),\qquad
\deg(f)=d\ \Longrightarrow\ N_K(f)\leq B_K(d).
\]
Determine whether the answers depend on \(K\). In particular, distinguish a bound allowed to depend on \(K\) from a single function \(B\colon\mathbb{Z}_{\geq 2}\to\mathbb{Z}_{\geq 0}\) satisfying
\[
N_K(f)\leq B(\deg f)
\]
for every admissible \(K\) and every \(f\in K(z)\) of degree at least two. The supplied excerpt does not resolve which field-uniform interpretation of the requested bound is intended.
\subsection{Short English statement}
For a rational map of degree at least two over the source's field \(\mathbb{C}_v\), ask whether there can be infinitely many periodic cycles of indifferent components that are not disks. If there are always finitely many, ask whether their number admits a bound depending only on the degree, and how the answers depend on the field. The field class and analytic terminology require clarification from the source.
\subsection{Sources}
\begin{itemize}
\item Local UnsolvedMath registry: permanent ID \(800412\), exact source problem code \texttt{catalogue entry}, title ``Open Problems in Arithmetic Dynamics --- Question 35.'' The supplied registry labels the problem \texttt{open}; this is supplied status metadata, not an independently verified mathematical or literature assessment.
\item Benedetto, DeMarco, Ingram, Jones et al., \emph{Current Trends and Open Problems in Arithmetic Dynamics} (2018), Question 35. Supplied source URL: \url{https://arxiv.org/abs/1806.04980}. Attribution and year are reproduced as supplied.
\item The source record reports access on 2026-10-16 and extraction by \texttt{arxiv\_latex}. These are supplied provenance details; no independent source visit or verification is claimed.
\end{itemize}
\end{document}
